% ── Style global du rapport UML-to-Code ──

\usepackage{placeins}
\usepackage{microtype}
\usepackage{enumitem}
\usepackage{xurl}
\usepackage{footnote}
\usepackage{etoolbox}

% Police : serif lisible (académique), pas sans-serif
\usepackage{lmodern}
\renewcommand{\familydefault}{\rmdefault}

% Espacement compact
\setstretch{1.12}
\setlength{\parskip}{0.35ex plus 0.15ex minus 0.1ex}
\setlist[itemize]{noitemsep, topsep=3pt, parsep=1pt, leftmargin=1.5em}
\setlist[enumerate]{noitemsep, topsep=3pt, parsep=1pt, leftmargin=1.8em}

% Liens et URL (coupure de ligne)
\hypersetup{
    colorlinks=true,
    linkcolor=black!80,
    filecolor=blue!70,
    urlcolor=blue!70!black,
    citecolor=black!70,
    breaklinks=true
}
\urlstyle{same}
\newcommand{\reporturl}[1]{\href{#1}{\nolinkurl{#1}}}

% Notes de bas de page
\makesavenoteenv{itemize}
\makesavenoteenv{enumerate}

% Tableaux améliorés
\newcolumntype{P}[1]{>{\raggedright\arraybackslash}p{#1}}
\newcolumntype{C}[1]{>{\centering\arraybackslash}p{#1}}
\renewcommand{\arraystretch}{1.12}
\setlength{\tabcolsep}{5pt}

\newcommand{\tablerule}{\midrule}
\newcommand{\tableheader}{\rowcolor{blue!12}\bfseries}

% En-têtes et pieds de page (référence projet en haut à droite)
\pagestyle{fancy}
\fancyhf{}
\fancyhead[L]{\small\textcolor{gray!70}{\leftmark}}
\fancyhead[R]{\small\textcolor{gray!65}{%
  Réf. : \href{https://github.com/ABAKAR5/uml-to-code}{GitHub}%
  \textbar\ \href{https://console.groq.com/docs}{Groq}%
  \textbar\ \href{https://uml-to-code-peach.vercel.app}{Vercel}%
}}
\fancyfoot[L]{\small\textcolor{gray!60}{INSTA — \annee}}
\fancyfoot[C]{\small\textbf{\thepage}}
\fancyfoot[R]{\small\textcolor{gray!60}{ABVIP}}
\renewcommand{\headrulewidth}{0.35pt}
\renewcommand{\footrulewidth}{0.25pt}
\renewcommand{\headrule}{\hbox to\headwidth{\color{blue!40}\leaders\hrule height \headrulewidth\hfill}}

% Pages préliminaires (sans en-tête gauche verbeux)
\fancypagestyle{prelim}{%
  \fancyhf{}
  \fancyfoot[C]{\small\textbf{\thepage}}
  \renewcommand{\headrulewidth}{0pt}
  \renewcommand{\footrulewidth}{0pt}
}

% Titres plus compacts
\titleformat{\chapter}[display]
  {\normalfont\huge\bfseries\color{blue!50!black}\centering}
  {\chaptertitlename\ \thechapter}
  {12pt}
  {\Huge}
\titlespacing*{\chapter}{0pt}{-18pt}{22pt}

\titleformat{\section}
  {\normalfont\Large\bfseries\color{blue!45!black}}
  {\thesection}
  {0.8em}
  {}
\titlespacing*{\section}{0pt}{10pt}{6pt}

\titleformat{\subsection}
  {\normalfont\large\bfseries}
  {\thesubsection}
  {0.6em}
  {}
\titlespacing*{\subsection}{0pt}{8pt}{4pt}

% Listings code plus compacts
\lstset{
    basicstyle=\ttfamily\footnotesize,
    aboveskip=6pt,
    belowskip=6pt
}

% Figures : légende serrée
\captionsetup{font=small, labelfont=bf, skip=6pt}

% Tableaux : taille réduite, placement souple (évite pages vides en bas)
\AtBeginEnvironment{table}{\small\centering}
\floatplacement{table}{htbp}
\floatplacement{figure}{htbp}
\raggedbottom

% TikZ (cycle de vie)
\usetikzlibrary{positioning,shapes,arrows.meta}
